% Options for packages loaded elsewhere
\PassOptionsToPackage{unicode}{hyperref}
\PassOptionsToPackage{hyphens}{url}
%
\documentclass[
  11pt,
]{article}
\usepackage{amsmath,amssymb}
\usepackage{iftex}
\ifPDFTeX
  \usepackage[T1]{fontenc}
  \usepackage[utf8]{inputenc}
  \usepackage{textcomp} % provide euro and other symbols
\else % if luatex or xetex
  \usepackage{unicode-math} % this also loads fontspec
  \defaultfontfeatures{Scale=MatchLowercase}
  \defaultfontfeatures[\rmfamily]{Ligatures=TeX,Scale=1}
\fi
\usepackage{lmodern}
\ifPDFTeX\else
  % xetex/luatex font selection
\fi
% Use upquote if available, for straight quotes in verbatim environments
\IfFileExists{upquote.sty}{\usepackage{upquote}}{}
\IfFileExists{microtype.sty}{% use microtype if available
  \usepackage[]{microtype}
  \UseMicrotypeSet[protrusion]{basicmath} % disable protrusion for tt fonts
}{}
\makeatletter
\@ifundefined{KOMAClassName}{% if non-KOMA class
  \IfFileExists{parskip.sty}{%
    \usepackage{parskip}
  }{% else
    \setlength{\parindent}{0pt}
    \setlength{\parskip}{6pt plus 2pt minus 1pt}}
}{% if KOMA class
  \KOMAoptions{parskip=half}}
\makeatother
\usepackage{xcolor}
\usepackage[margin=25mm]{geometry}
\setlength{\emergencystretch}{3em} % prevent overfull lines
\providecommand{\tightlist}{%
  \setlength{\itemsep}{0pt}\setlength{\parskip}{0pt}}
\setcounter{secnumdepth}{-\maxdimen} % remove section numbering
\usepackage{amsmath,amssymb,mathtools,xurl}
\emergencystretch=3em
\ifLuaTeX
  \usepackage{selnolig}  % disable illegal ligatures
\fi
\usepackage{bookmark}
\IfFileExists{xurl.sty}{\usepackage{xurl}}{} % add URL line breaks if available
\urlstyle{same}
\hypersetup{
  pdftitle={Positive-measure avoidance with growing logarithmic gaps},
  pdfauthor={mxym/math research project (AI-assisted)},
  hidelinks,
  pdfcreator={LaTeX via pandoc}}

\title{Positive-measure avoidance with growing logarithmic gaps}
\author{mxym/math research project (AI-assisted)}
\date{8 October 2026}

\begin{document}
\maketitle

\subsection{Abstract}\label{abstract}

We prove positive-measure affine non-universality for decreasing null
sequences whose consecutive logarithmic gaps are \(o(\log\log z)\),
where \(z=-\log_2 a\). This includes sequences whose adjacent ratios
tend to zero and whose occupied logarithmic bins have upper Banach
density zero. A more general theorem needs this filling condition only
on arbitrarily late multiplicative logarithmic windows. For each
prescribed countable family we construct a closed periodic set of
measure arbitrarily close to one that misses infinitely many distinct
outputs in every tail of every nonzero leading-power map with a
power-controlled remainder. The proof allows the branching, depth and
separation gap of a finite routing tree to grow with its position. Its
new schedule has total span \(U^{o(1)}\), which leaves sufficient room
for the error buffer.

The finite routing construction is inherited and explicitly reproved
below. The contributions here are the variable-parameter schedule,
local-window sampling, the resulting zero-density sequence classes and
the intermittent window extension. No complete solution of the Erdős
similarity conjecture, world-first claim or external human peer review
is asserted.

\subsection{1. Problem, notation and
statements}\label{problem-notation-and-statements}

All unmarked logarithms are natural. Lebesgue measure is denoted by
\(m\). For a measurable one-periodic set \(S\), put
\(\operatorname{dens}(S)=m(S\cap[0,1))\). A positive null sequence is
strictly decreasing, \(a_n>0\), \(a_n\to 0\). Discarding finitely many
initial terms will always be permitted when selecting test inputs, and
never weakens the eventual avoidance conclusion. Put
\(z_n=-\log_2 a_n\); eventually \(z_n>e\), and \(z_n\) increases to
infinity.

An affine copy of \(A=\{a_n\}\) is \(y+cA\) with \(c\ne 0\). The Erdős
similarity conjecture asks whether every infinite subset of the line is
absent in this sense from some positive-measure measurable set. We treat
a specified class of null sequences, not all infinite sets.

\textbf{Definition 1 (filling on late annuli).} A sequence \((z_n)\) has
property \(W\) if, for every integer \(R\ge 2\), there are integers
\(U_j\to \infty\) and real numbers \(D_j\ge 1\) with

\[
\frac{D_j}{\log\log U_j}\longrightarrow0,                 \tag{1}
\]

such that every closed interval \([v,v+D_j]\) contained in
\([U_j/R,RU_j]\) contains a member of \((z_n)\). Only sufficiently large
\(j\) are relevant, so all logarithms in (1) have positive arguments.

\textbf{Theorem 2 (local-window avoidance).} Let \((A_l)\) be a
prescribed nonempty countable family of positive null sequences whose
logarithms have property \(W\). For every \(0<\varepsilon <1\) there is
a closed nowhere dense one-periodic set \(E\subset\mathbb R\) with
\(m(E\cap [t,t+1])>1-\varepsilon\) for every real \(t\), with the
following property. After \(E\) has been chosen, fix any family member
and any \(s,\alpha >0\), \(c\ne 0\), \(y\in\mathbb R\), and \(M\ge 0\).
If \(f\) is defined on a sufficiently small tail of that member and

\[
 |f(a)-y-ca^s|\le M a^{s+\alpha}                           \tag{2}
\]

there, then in every positive input tail there are infinitely many
distinct values of \(f(a)\) outside \(E\). In particular
\(F=E\cap [0,1]\) is compact, \(m(F)>1-\varepsilon\), and contains no
nontrivial affine copy of any family member.

The family is fixed before \(E\); the theorem does not construct a
single set for all null sequences. There are no regularity or
injectivity assumptions on \(f\) beyond (2). Its domain need only be the
stated tail.

\textbf{Corollary 3 (growing consecutive gaps).} Theorem 2 applies
whenever

\[
 z_{n+1}-z_n=o(\log\log z_n).                              \tag{3}
\]

\textbf{Corollary 4 (ratios tending to zero).} For every fixed
\(0<\beta <1\) it applies to

\[
 a_n=2^{-n(\log\log(n+20))^\beta},\qquad n\ge1.            \tag{4}
\]

For this sequence \(a_{n+1}/a_n\to 0\), and its occupied dyadic
logarithmic bins have upper Banach density zero. One may also prescribe
countably many such sequences at once. The assertion is not simultaneous
over an unspecified uncountable family of \(\beta\).

\subsection{2. Sampling from a finite logarithmic
annulus}\label{sampling-from-a-finite-logarithmic-annulus}

Fix one sequence, a requested input tail \(a<2^{-h}\), and a compact
exponent rectangle \(0<s_0\le s_1\). Fix an integer coefficient scale
\(k\). Write

\[
 \delta=4/s_0,\qquad R\ge\max(2,4s_1,8/s_0),              \tag{5}
\]

with integer \(R\). Choose the annuli from Definition 1 for this \(R\).
All subsequent choices are made separately on the \(j\)th annulus.

Start with a logarithm in \([U_j/R,U_j/R+D_j]\), which exists for large
\(j\). Thereafter select greedily the first original logarithm strictly
larger than the previous selected one plus \(\delta\). Continue
indefinitely, not just inside the annulus. Denote the selected
logarithms by \(\zeta _i\) and the inputs by \(b_i=2^{-\zeta _i}\). For
large \(j\) every selected input is in the requested tail. Consecutive
selected logarithms have gaps greater than \(\delta\).

Whenever \(\zeta _i+\delta +1/2+D_j\le RU_j\), the interval
\([\zeta _i+\delta +1/2,\zeta _i+\delta +1/2+D_j]\) contains an original
logarithm. The first eligible logarithm therefore satisfies

\[
 \delta<\zeta_{i+1}-\zeta_i<\delta+D_j+1.                 \tag{6}
\]

Put

\[
 D=s_1(\delta+D_j+2),\quad \eta=1/(2D).                    \tag{7}
\]

In particular \(D\ge 4\), and \(D=o(\log\log  U_j)\) along these annuli.
Suppose a routing template lies in \([U,U+T]\), with \(U=U_j\) and
\(T\le U\). For all sufficiently large \(j\), its entire relevant
logarithm range, including the extra \(\delta +D_j+1\), lies strictly
below \(RU\). The first selected logarithm lies below \((U+k)/s_1\);
these assertions follow from (5), \(D_j=o(U)\), and the fact that
\(k,\delta\) are fixed.

For each \(s\in [s_0,s_1]\), set \(w_i=s\zeta _i-k\). The sequence
starts below \(U\), is separated by more than three, and every gap
needed to cross a window inside \([U,U+T]\) is less than \(D\), by (6).
Thus for every open window \((u,u+\ell )\) in that template, if
\(\ell \ge 2D\),

\[
 \eta\ell\le\#\{i:u<w_i<u+\ell\}\le1+\ell/3.            \tag{8}
\]

For the lower bound, the last member at or below \(u\) and the first
member at or above \(u+\ell\) enclose \(r\) strictly interior members.
The length is at most \((r+1)D\), so \(r\ge \ell /D-1\ge \ell /(2D)\).
This argument includes equality at either endpoint. The upper bound uses
separation. None of (6)--(8) requires a bound on the logarithmic gaps
after leaving the annulus.

\subsection{3. A uniform finite routing
lemma}\label{a-uniform-finite-routing-lemma}

Here we give the finite construction with all its numerical guards. They
are valid for each parameter choice separately; their proof never
requires the tree parameters to be fixed while \(U\) tends to infinity.

Fix \(0<p<1/12\). Choose integers \(b\ge 2\), \(d\ge 1\), \(g\ge 1\),
and \(L\ge 2\). Use a complete ordered \(b\)-ary tree of depth \(d\),
with \(K=\sum _{r=1}^d b^r\) edges. Recursively define

\[
 T_1=bL+(b-1)g,\qquad
 L_r=g+T_{r-1},\quad
 T_r=2bT_{r-1}+(3b-1)g\quad(r\ge2).                       \tag{9}
\]

Place an integer window of length \(L_r\) on each edge from a
height-\(r\) vertex. The edge is followed by its whole child subtree,
and consecutive pieces are separated by \(g\) unused integers. The
template starts at \(U\). Its total span is \(T=T_d\). An edge window of
length \(\ell =L_r\), together with its subtree, has span at most
\(2\ell\). All lengths are at least \(L\). These facts follow directly
from (9), including the height-one case. More explicitly, the preorder
traversal places \(g\) unused indices between every consecutive pair of
edge windows. A window starting at \(u\) ends at \(v=u+\ell -1\); the
last endpoint of the whole template is \(U+T-1\).

Assume the local samples obey (8), and the following guards hold:

\[
\begin{gathered}
 U\ge4,\quad U>|k|,\quad T\le U,\quad L\ge2D,\quad L\ge g,\qquad
 \lambda=p\eta(b-1)/2,\quad \kappa=\lambda-4\log2>0,\\
 (1-2^{1-b})^d<p,\qquad K2^{3-g}<p,\\
 46080b^2(U+1)^2e^{-\kappa L}<p.                           \tag{10}
\end{gathered}
\]

The initially selected sample must precede \(U\) in every output
exponent, as proved in Section 2. Finally fix an error exponent
\(\alpha _0>0\) and an integer error bound \(q\ge 1\). Let

\[
 N=2^{U+T+2},\quad
 r_U=(q2^{-k}+1)2^{-U}2^{-\alpha_0(U+k)/s_1},\quad
 4Nr_U<p.                                                 \tag{11}
\]

\textbf{Lemma 5 (finite blocker).} Under these assumptions there is an
open one-periodic set \(H\) of density less than \(6p\) such that, for
every \(x\in\mathbb R\), \(s\in [s_0,s_1]\), \(t\in [1,2]\) and
arbitrary errors \(|e(b_i)|\le q b_i^{s+\alpha _0}\), at least one
selected input satisfies

\[
 x+t2^k b_i^s+e(b_i)\in H.                                 \tag{12}
\]

The index in (12) may lie beyond the finite template. This is essential
for repairing exceptional centers.

\subsubsection{3.1 Tables, stable centers and actual
addresses}\label{tables-stable-centers-and-actual-addresses}

At integer endpoint \(v\) use the periodic grid of \(2^{v+3}\) half-open
cells. The address of a real point is
\(\lfloor2^{v+3}x\rfloor\bmod 2^{v+3}\), and cells are left closed and
right open. At each nonlast outgoing edge put independent fair bits
indexed by its endpoint-grid cells. At each leaf put independent
Bernoulli-\(p\) bits indexed by its incoming endpoint grid. Entries of
different tables are independent. Route a point to the first nonlast
child whose bit is one, or to the last child when all those bits are
zero. Include the point in \(B\) exactly when its reached terminal bit
is one. The set \(B\) is a union of cells on the finest \(N\)-grid. For
a fixed point, conditioning on the selector tables leaves its terminal
entry Bernoulli-\(p\); hence \(\mathbb E\,\operatorname{dens}(B)=p\).

If a window begins at \(u=v+g+1\) after a preceding endpoint \(v\), call
a center unstable when a boundary of the \(2^{v+3}\) grid lies in
\((x,x+2^{-v-g}]\). Over one period its measure is at most \(2^{3-g}\).
The unstable union has measure at most \(K2^{3-g}<p\).

An active offset on a length-\(\ell\) window obeys

\[
 2^{-u-\ell}<t2^{k-s\zeta_i}<2^{1-u}.                     \tag{13}
\]

At a stable center it shares all earlier grid addresses with the center,
since (13) is smaller than the stability allowance. At its own grid its
address differs from the center's; the grid width is one quarter of the
lower bound in (13). Different active inputs have offsets in a ratio
greater than eight, so their addresses differ too. Offsets lie in an
interval of length less than \(1/8\); periodic wrapping cannot defeat
these separations. The distinctions persist on every finer grid. Exact
grid boundaries use the stated half-open convention.

\subsubsection{3.2 Conditional probability at a first
default}\label{conditional-probability-at-a-first-default}

At a stable center expose its addressed entry in every selector table.
This determines its entire route. The probability that its route has no
default child is exactly \((1-2^{1-b})^d<p\). Otherwise condition on the
complete exposure atom and its first default, at height \(r\).

For each of the \(b-1\) nondefault edges from that vertex and each
active input, test its own selector bit and then the terminal bit found
by routing it in that child subtree. Earlier addresses agree with the
center; the required preceding nondefault bits are zero. A successful
test is therefore a point of the actual global set \(B\), not of a
hypothetical local set.

There are at least \(m_*\ge \eta (b-1)L_r\) tests. Their own selector
entries are distinct and were not exposed. Condition next on all
selectors. Points that reach different leaves use different terminal
tables. Points that reach the same leaf have distinct addresses already
on their common own edge, and hence distinct finer terminal addresses.
Consequently the conditional probability of failure of all tests is
\((1-p)^{\sum  A_i}\), where \(A_i\) are the tested own selector bits.
Averaging the independent fair \(A_i\) gives

\[
 (1-p/2)^{m_*}\le e^{-\lambda L_r}.                        \tag{14}
\]

Shared auxiliary selector reads do not affect this calculation. Entire
routing paths are not being assumed independent. Multiplying (14) by the
exposure-atom mass yields a joint estimate, also valid for zero-mass
atoms.

\subsubsection{3.3 The continuum of parameters and all boundary
strata}\label{the-continuum-of-parameters-and-all-boundary-strata}

Fix the center and the first-default vertex, and put \(\ell =L_r\).
There are at most

\[
 P=(b-1)\left(\left\lceil (U+T+|k|)/3\right\rceil+1\right) \tag{15}
\]

candidate edge--input pairs: any candidate logarithm is positive, is
below \((U+T+|k|)/s_0\), and is separated from the others by more than
\(3/s_0\). For each pair retain the two activation equations. Retain
also all boundaries of its local finest grid whose positive offsets from
the center are less than \(2^{1-u}\). Its local span is at most
\(2\ell\), so there are at most \(1+2^{2\ell +3}\) such boundaries.

Writing \(v=\log  t\), a boundary of positive offset \(\beta\) is
crossed at

\[
 v-s\zeta_i\log2+k\log2-\log\beta=0.                     \tag{16}
\]

All retained equations are affine lines in \((s,v)\). Their full sign
vectors determine activation and each actual local finest-grid address.
Inactive values at activation equality and right-cell floor values at
grid equality are encoded by zero signs. At most \(m=P(3+2^{2\ell +3})\)
lines occur.

For \(m\) affine lines in the plane the total number of distinct sign
vectors is at most \(2m^2+1\): after merging coincident lines, there are
at most \(m(m-1)/2\) vertices, \(m^2\) open line edges, and
\(1+m(m+1)/2\) faces. Each stratum has a constant full sign vector. This
count also handles parallel lines, repeated intersections and constant
equations. Intersect the strata with the closed parameter rectangle and
select one point from each realized sign vector. No table entries are
used in that choice.

The resulting representative count is bounded by \(20(m+5)^2\). Union
bounding (14) over these representatives gives a continuum miss
probability on the exposure atom at most its mass times

\[
 5120b^2(U+T+|k|+2)^2e^{-\kappa\ell}
 \le46080b^2(U+1)^2e^{-\kappa L}<p.                       \tag{17}
\]

Indeed \(P\le b(U+T+|k|+2)\) and, for \(P\ge 1,\ell \ge 2\),
\(P(3+2^{2\ell +3})+5\le 16P2^{2\ell }\). Squaring accounts for
\(4\log 2\) in \(\kappa\); \(T\le U\) and \(|k|<U\) account for the
factor nine in (17). The representatives work for every assignment of
unexposed bits. They need not be selected measurably as functions of the
center; only a pointwise bound is used before a finite sum and
integration.

\subsubsection{3.4 Closed exceptional centers and open
repair}\label{closed-exceptional-centers-and-open-repair}

Let \(B_1=B+(-r_U,r_U)\) and \(B_2=B+(-2r_U,2r_U)\). Enlarging a union
of finest-grid cells by radius \(2r_U\) increases its periodic density
by at most \(4Nr_U\). Therefore
\(\mathbb E\,\operatorname{dens}(B_2)<2p\). There are finitely many
potentially active inputs: their logarithms are bounded above by
\((U+T+|k|)/s_0\). Let \(I\) contain them, and let \(S_i\) be the
relatively open set of exponents for which \(s\zeta _i-k\) belongs to
the union of all the open template windows. Define \(Q\) as the set of
centers for which some parameter pair has, for every \(i\in I\), either
\(s\notin S_i\) or its unperturbed output outside \(B_1\).

Each inactivity condition is closed in the closed parameter rectangle,
because activation windows are open. Each outside-\(B_1\) condition is
closed by continuity. Their finite intersection is closed, and
projection along the compact rectangle is closed. Thus \(Q\) is closed
and periodic. At stable centers, membership forces a continuum miss at
the first default or a route with no default. Equations (10) and (17)
bound its probability by \(2p\). Integrating and allowing the unstable
union gives \(\mathbb E\,\operatorname{dens}(Q)\le 3p\). This uses only
finite sums and measurable indicators.

Choose an outcome with
\(\operatorname{dens}(B_2)+\operatorname{dens}(Q)\le 5p\). Thicken the
closed periodic \(Q\) to an open periodic \(V\) with
\(\operatorname{dens}(V)<\operatorname{dens}(Q)+p\); continuity from
above on a finite period justifies this, including the empty-set case.
Put \(H=B_2\cup V\). Its density is less than \(6p\).

If \(x\notin Q\), some active input has unperturbed output in \(B_1\).
Activation gives \(s\zeta _i-k>U\), so its allowed error satisfies

\[
 q b_i^{s+\alpha_0}
 \le q2^{-k}2^{-U}2^{-\alpha_0(U+k)/s_1}<r_U.              \tag{18}
\]

That perturbed output belongs to \(B_2\). If \(x\in Q\subset V\), use
the full selected tail: the perturbed offsets converge to zero uniformly
under the fixed bounds \(2^{k+1}b_i^{s_0}+q b_i^{s_0+\alpha _0}\). An
output is eventually in the open neighborhood \(V\) of \(x\). This input
need not be active in the finite template. These two cases prove Lemma
5.

\subsection{4. Growing trees fit on the same
annulus}\label{growing-trees-fit-on-the-same-annulus}

We now prove the main new scheduling step. In this section all tree
parameters are allowed to depend on \(j\), equivalently on \(U=U_j\).
They are chosen before any table outcome or center.

Take \(D\) from (7), and set

\begin{equation*}
\begin{aligned}
 b&=2+\lceil24D/p\rceil,\qquad
 d=\left\lceil2^{b-1}\log(2/p)\right\rceil,\\
 K&=\sum_{r=1}^d b^r,\qquad
 g=1+\left\lceil\log_2(8K/p)\right\rceil,\\
 \lambda&=p(b-1)/(4D),\qquad \kappa=\lambda-4\log2,\\
 L&=\max\left(g,\lceil2D\rceil,2,
 \left\lceil\frac{4\log(U+2)+\log(46080b^2/p)+1}{\kappa}\right\rceil\right).
\end{aligned}
                                                               \tag{19}
\end{equation*}

These are finite integers, however enormous. Since \(\lambda \ge 6\) and
\(\log 2<3/4\), \(\kappa >3\). Also

\[
 (1-2^{1-b})^d\le e^{-2^{1-b}d}\le p/2<p,\qquad
 K2^{3-g}<p.                                                  \tag{20}
\]

The definition of \(L\) gives the last entropy guard directly:

\[
46080b^2(U+1)^2e^{-\kappa L}
\le pe^{-1}(U+1)^2/(U+2)^4<p.                                \tag{21}
\]

\textbf{Lemma 6 (uniform tree-span bound).} For every real
\(b\ge 2,g,L\ge 0\) and integer \(r\ge 1\), the recurrence (9) satisfies

\[
 T_r+2g\le b(2b)^{r-1}(L+2g),\qquad
 T_r\le(2b)^r(L+2g).                                        \tag{22}
\]

\emph{Proof.} At \(r=1\), \(T_1+2g=bL+(b+1)g\le b(L+2g)\). For
\(r\ge 2\), \(T_r+2g=2bT_{r-1}+(3b+1)g\le 2b(T_{r-1}+2g)\). Induction
proves the first inequality. The second follows from \(g\ge 0\),
\(L+2g\ge 0\), and \(b\le 2b\). This proof is uniform in depth and
branching. \(\square\)

\textbf{Lemma 7 (sublinear span).} With (19) and \(D=o(\log\log  U)\),

\[
 L=O_p(\log U),\qquad T=T_d=U^{o(1)},\qquad T/U\to0.        \tag{23}
\]

\emph{Proof.} Put \(t=\log  U\), so \(D=o(\log  t)\) and
\(b=o(\log  t)\). For every fixed \(\theta >0\),

\[
 d\log(2b)=t^{o(1)}=o(t).                                   \tag{24}
\]

Here is a direct justification, also covering bounded \(b\). We have
\(d\le 1+2^{b-1}\log (2/p)\). The logarithm of the right-hand side times
\(\log (2b)\) is at most \(b \log 2+O_p(1)+\log\log (2b)=o(\log  t)\).
Thus the expression is eventually at most \(t^\theta\); take
\(\theta <1\). The geometric sum gives \(K\le 2b^d\), and hence
\(g=O_p(1+d \log (2b))=t^{o(1)}=o(t)\). The lower-bound terms \(g,2D,2\)
in (19) are \(o(t)\) or bounded. The numerator of its final term is
\(4t+o(t)+O_p(1)\) and \(\kappa >3\), so \(L=O_p(t)\). By (22),

\[
 \log T\le d\log(2b)+\log(L+2g)=o(t).
\]

Consequently for every \(\theta >0\), \(T\le U^\theta\) eventually. In
particular \(T/U\to 0\). Unlike a fixed-tree asymptotic argument, this
calculation retains the changing depth and all changing coefficients.
\(\square\)

For large \(j\) we have \(T\le U\), all the annulus margins in Section
2, and every standing integer guard. These statements follow along the
very same chosen annuli; no independent choice of a later \(U\) with an
older gap bound is made. Finally (11) obeys

\[
 4Nr_U=2^{T+4}(q2^{-k}+1)2^{-\alpha_0(U+k)/s_1}\to0,        \tag{25}
\]

because its logarithm is \(-(\alpha _0 \log 2/s_1)U+o(U)+O(1)\). Thus
one sufficiently late annulus supplies all guards of Lemma 5. This
proves a blocker of density less than \(6p\) for every prescribed tail,
bounded exponent rectangle, coefficient scale and error bound.

\subsection{5. Completion of the global
theorem}\label{completion-of-the-global-theorem}

Enumerate the countable tuples \((l,N,j,k,q,h)\), where \(N\ge 2\),
\(j,q,h\ge 1\) are integers and \(k\in\mathbb Z\). Assign strictly
positive budgets \(p_\tau <1/12\) with
\(12\sum _\tau  p_\tau <\varepsilon\). Apply the blocker just proved
with \([s_0,s_1]=[1/N,N]\), \(\alpha _0=1/j\), the requested sequence
and input tail, and coefficient scale \(k\). Let its open periodic set
be \(H_\tau\). Define

\[
 \mathcal U=\bigcup_\tau(H_\tau\cup(-H_\tau)),\qquad
 E=\mathbb R\setminus\mathcal U.                             \tag{26}
\]

Then \(E\) is closed and periodic and its measure in a period is greater
than \(1-\varepsilon\). For a measurable one-periodic set, splitting any
unit interval at integers and translating its two pieces shows the same
measure for every unit interval. Endpoints have measure zero.

Fix the data in (2). Choose \(N\) with \(s\in [1/N,N]\), \(j\) with
\(1/j\le \alpha\), and integer \(q\ge \max (1,M)\). For \(a<1\) its
remainder is at most \(q a^{s+1/j}\). When \(c>0\), write \(c=2^k t\),
\(1\le t<2\). For every sufficiently large \(h\), the selected inputs
are inside the actual domain of \(f\) and the eventual remainder bound.
Lemma 5 at the tuple with that \(h\) produces a value in
\(H_\tau \subset\mathcal U\). When \(c<0\), apply it to \(-f\) and then
reflect; the value belongs to \(-H_\tau\).

We therefore obtain missed values at inputs tending to zero. Eventually
\((1/2)|c|a^s\le |f(a)-y|\le (3/2)|c|a^s\), so these values are
different from \(y\) and tend to it. A finite set avoiding \(y\) cannot
have that property. They give infinitely many distinct missed values
inside every input tail. No injectivity was used. If \(E\) had an open
interval, take an interior center and the map \(f(a)=y+a\) on a small
tail of any one family member; all outputs would be in that interval, a
contradiction. Nonemptiness of the family is used here. Thus \(E\) is
nowhere dense. Intersecting with \([0,1]\) proves the compact assertion.
This finishes Theorem 2. \(\square\)

\subsection{6. Consecutive gaps, explicit zero-density examples, and
intermittent
windows}\label{consecutive-gaps-explicit-zero-density-examples-and-intermittent-windows}

\emph{Proof of Corollary 3.} Discard a finite prefix so the logarithms
are positive. Define \(\Omega (Z)=\max \{z_{n+1}-z_n:z_n\le Z\}\), using
zero for an empty maximum. There are only finitely many entries under
any finite \(Z\). Hypothesis (3) implies \(\Omega (Z)=o(\log\log  Z)\):
for any positive \(\eta\), all sufficiently late entries are at most
\(\eta  \log\log  z_n\le \eta  \log\log  Z\); the finitely many early
entries contribute a fixed bound, whose ratio tends to zero. For a fixed
integer \(R\), set \(U_j=j\) and \(D_j=1+\Omega (Rj)\) once \(j\) is
large. In any closed interval \([v,v+D_j]\subset [j/R,Rj]\), the last
logarithm at or below \(v\) is followed by one at distance at most
\(\Omega (Rj)<D_j\) from it. Unless the left endpoint already is a
logarithm, this successor is strictly inside the interval. Thus
Definition 1 holds. Since \(\log\log (Rj)/\log\log  j\to 1\), its ratio
condition holds too. \(\square\)

\emph{Proof of Corollary 4.} Put \(\ell (x)=\log\log (x+20)\) and
\(z(x)=x \ell (x)^\beta\). Then

\[
 z'(x)=\ell(x)^\beta+
 \frac{\beta x\ell(x)^{\beta-1}}{(x+20)\log(x+20)}
       \sim(\log\log x)^\beta.                              \tag{27}
\]

The mean-value theorem gives \(z_{n+1}-z_n\sim (\log\log  n)^\beta\),
while \(\log\log  z_n\sim \log\log  n\). Their ratio tends to zero
because \(\beta <1\); the gap itself tends to infinity because
\(\beta >0\). Thus adjacent ratios tend to zero. The integer parts of
the \(z_n\) are eventually distinct and have gaps tending to infinity.
For each integer \(M\), discard finitely many bins until these gaps are
at least \(M\). Every interval of \(L\) consecutive bins then contains
at most \(C_M+1+L/M\) occupied bins. Taking its supremum, then
\(\limsup\) divided by \(L\), and then \(M\to \infty\), proves upper
Banach density zero. The existing positive-density criterion does not
apply. \(\square\)

Property \(W\) is strictly broader than (3). Here is an explicit
example. For integers \(j\ge 2\), let \(U_j=2^{2^{j^2}}\) and take all
logarithms

\[
 Z_j=\{\lceil U_j/j\rceil+mj:m\ge0,
                    \ \lceil U_j/j\rceil+mj\le jU_j\}.     \tag{28}
\]

List the disjoint blocks \(Z_j\) in increasing order and set
\(a_n=2^{-z_n}\). For each fixed \(R\), all sufficiently large \(j>2R\)
have \([U_j/R,RU_j]\) strictly inside the span of \(Z_j\), with margins
greater than \(2j\). Every length-\(2j\) subinterval there contains a
lattice point. Also \(2j/\log\log  U_j\to 0\), since
\(\log\log  U_j=j^2 \log 2+\log\log 2\). Thus property \(W\) holds. The
gap between successive blocks is far larger than \(\log\log  z_n\), so
(3) fails. Within a late block gaps equal \(j\), and between blocks they
are larger; the same bin-counting argument proves zero upper Banach
density. Theorem 2 therefore also treats sequences with arbitrarily
large gaps outside their useful annuli.

\subsection{7. Prior work, proof dependence and
limitations}\label{prior-work-proof-dependence-and-limitations}

The finite center-exposure routing, parameter sign representatives and
open repair used in Section 3 were developed in the geometric/continuum
power-avoidance work publicly recorded in \texttt{mxym/math}, itself
using the geometric-case method of OpenAI Mathematics. The earlier
theorem assumed globally bounded occupied logarithmic gaps. Its complete
release was replayed from all 65 owned source modules, with 568 public
theorems and 35,620 safe-closure declarations checked in empty Lean
kernel environments. The relevant published files and hashes are listed
in \texttt{SOURCES.json}. We do not substitute that formalized
bounded-gap endpoint for the new claim: Sections 2 and 4 are the added
arguments and have complete written proofs; the whole new endpoint is
not yet Lean-formalized.

Cruz--Lai--Pramanik, \emph{Large Sets Avoiding Affine Copies of Infinite
Sequences}, Real Analysis Exchange 48(2) (2023), 251--270,
arXiv:2204.12720, construct a dimension-one Borel avoiding set under a
slow-growth condition on \(\log\log (a_n/a_{n+1})\). Its positive
measure is not established; Section 4.4 explicitly discusses this
missing point. Their dimension-one subset \(O\) has measure zero; this
does not mean that their whole avoiding set \(E_K\) is proved to have
measure zero. Our specified examples have positive-measure avoiding
sets, indeed measure arbitrarily close to one. This is a distinction in
the mathematical conclusions, not a certified priority assertion based
on an exhaustive literature survey.

Feng--Lai--Xiong, arXiv:2312.01319, prove bi-Lipschitz universality for
suitable lacunary sequences, including (4). That uses transformations
more general than affine maps and does not contradict this result. The
earlier geometric case and the repository's positive logarithmic-density
criterion do not cover the zero-density examples (4) and (28).

For \(a_n=2^{-n^2}\), consecutive logarithmic gaps are of order
\(\sqrt(z_n)\); for \(a_n=2^{-2^n}\), they are of order \(z_n\). Neither
satisfies our hypotheses, nor the local annular condition. Both, and
arbitrary null sequences, remain outside our result. The proof also does
not cover arbitrary flat germs, zero leading coefficients, or an
uncontrolled \(o(a^s)\) error. All sets are existence constructions; the
finite tables can be astronomically large. The exact checker supplies
algebraic and finite controls, not a computed high-measure avoiding set.
The whole analytic theorem is not represented as independently
human-reviewed or fully formalized.

\end{document}
